\section{Related Work}
\label{sec:related}

STaR and ReST$^{\mathrm{EM}}$ learn from filtered model-generated solutions
\citep{zelikman2022star,singh2023restem}. Recursive self-improvement also
encompasses evolving exploration policies, as in
Dream-RSI~\citep{zheng2026dreamrsi}. Existing theory studies finite-sample
reward improvement and curricula \citep{liu2026taskcentric} or solver--verifier
gaps \citep{sun2025solververifier}; we instead analyze incorrect targets at
matched acceptance rates. Systematic false positives can impair RLVR
\citep{egashira2026systematic}, while bandit phase boundaries
\citep{rad2026rate} and reward/loss corrections
\citep{cai2025noisy,natarajan2013noisy,patrini2017correction} depend on their
learner and noise assumptions. Our within-class intervention exposes gradient
structure left undetermined by aggregate rates, without contradicting those
guarantees. Influence functions and TracIn connect training examples to
predictions \citep{koh2017influence,pruthi2020tracin}; GRAD-MATCH and
PartitionSel constrain selected gradients
\citep{killamsetty2021gradmatch,agrawal2026partitionsel}. Our contribution is
non-identifiability at matched summaries, population-flow separation, and a
policy-specific audit threshold, not a new gradient-selection objective.
Diagnostics use actual optimizer displacements rather than inverse-Hessian
influence and do not certify task-risk improvement. Further comparisons
appear in Appendix~\ref{app:extended-related-work}.
